%%%%%%%%%%%%%%%%%%%%%%%%%%%%%%%%%%%%%%%%%%%%%%%%%%%%%%%%%%%%%%%%%%%%%%%%%%%%%%%
%%  TKS_Overton_Window_Manipulation.tex
%%  Overton Window Manipulation: Reality Tunneling
%%
%%  Comprehensive TKS Formalization of Discourse Boundary Dynamics
%%  Part of TKS v7.4 Scenario Equation System
%%  ADDITIVE to v7.4 canon
%%%%%%%%%%%%%%%%%%%%%%%%%%%%%%%%%%%%%%%%%%%%%%%%%%%%%%%%%%%%%%%%%%%%%%%%%%%%%%%

%% ============================================================================
%% CHAPTER: OVERTON WINDOW MANIPULATION
%% ============================================================================
\chapter{Overton Window Manipulation: Threat Model for Discourse Boundary Attacks}
\label{ch:overton-window}

\begin{tcolorbox}[colback=red!5!white,colframe=red!75!black,title=\textbf{THREAT MODEL NOTICE},breakable]
The following models are presented as \textbf{threat models for defenders}. They describe how hostile actors---authoritarian regimes, disinformation campaigns, or extremist movements---might systematically shift the boundaries of acceptable discourse to normalize previously unthinkable positions.

Understanding these mechanisms is essential for:
\begin{itemize}
    \item \textbf{Media literacy educators}: Teaching critical awareness of gradual normalization
    \item \textbf{Platform integrity teams}: Detecting coordinated window-shifting campaigns
    \item \textbf{Democratic institutions}: Building resilience against discourse manipulation
    \item \textbf{Researchers}: Developing early-warning systems for radicalization pipelines
\end{itemize}
The mathematics is value-neutral; the defensive framing ensures responsible documentation.
\end{tcolorbox}

\begin{tcolorbox}[colback=blue!5!white,colframe=blue!75!black,title=\textbf{Chapter Overview},breakable]
This chapter formalizes the \textbf{Overton Window}---the range of ideas considered
acceptable in public discourse at any given time---as a \textit{psychological attractor basin}
in idea space $\Ideas$. Using canonical TKS machinery (transfinite fractals, Collage Theorem,
ACBE descent), we model:
\begin{itemize}
    \item How acceptability boundaries form and stabilize
    \item Gradual boundary shifts via transfinite sequences
    \item The design of shift strategies under constraint optimization
    \item Normalization dynamics governing acceptance thresholds
\end{itemize}
The mathematics presented here is value-neutral: the same equations describe
civil rights expansion, technology adoption shifts, and deliberate manipulation campaigns.
\end{tcolorbox}

%% ============================================================================
%% SECTION: MATHEMATICAL DEFINITIONS
%% ============================================================================
\section{Mathematical Definitions}
\label{sec:overton-definitions}

\subsection{The Overton Window as Attractor Basin}

\begin{definition}[Discourse Space]
\label{def:discourse-space}
The \textbf{discourse space} $\mathcal{D} \subset \Ideas$ is the subspace of Ideas
pertaining to a specific domain of public discussion (politics, ethics, technology, etc.).
Elements of $\mathcal{D}$ are \textit{positions}---discrete stances or policy proposals
that can be articulated and debated.
\end{definition}

\begin{definition}[Overton Window]
\label{def:overton-window}
The \textbf{Overton Window} $\mathcal{O} \subset \mathcal{D}$ is the attractor basin
in discourse space representing the range of ideas considered ``acceptable'' by the
prevailing social consensus at time $t$. Formally:
\[
\boxed{
    \mathcal{O}(t) := \left\{ x \in \mathcal{D} \;\Big|\;
    \mu_{\mathrm{accept}}(x, t) \geq \theta_{\mathrm{norm}}(t) \right\}
}
\]
where:
\begin{itemize}
    \item $\mu_{\mathrm{accept}} : \mathcal{D} \times \mathbb{R}_{\geq 0} \to [0,1]$
          is the \textbf{acceptability measure} assigning each position a social
          acceptance probability at time $t$
    \item $\theta_{\mathrm{norm}}(t) \in [0,1]$ is the \textbf{normalization threshold}---
          the minimum acceptability required for inclusion in mainstream discourse
\end{itemize}
\end{definition}

\begin{definition}[Window Topology]
\label{def:window-topology}
The Overton Window has canonical boundary structure:
\begin{align}
    \partial^{-}\mathcal{O} &:= \left\{ x \in \mathcal{D} \;\Big|\;
        \mu_{\mathrm{accept}}(x) = \theta_{\mathrm{norm}} - \epsilon \right\}
        \quad \text{(unthinkable boundary)} \\
    \partial^{+}\mathcal{O} &:= \left\{ x \in \mathcal{D} \;\Big|\;
        \mu_{\mathrm{accept}}(x) = \theta_{\mathrm{norm}} + \epsilon \right\}
        \quad \text{(radical boundary)} \\
    \mathrm{Core}(\mathcal{O}) &:= \left\{ x \in \mathcal{D} \;\Big|\;
        \mu_{\mathrm{accept}}(x) \geq \theta_{\mathrm{policy}} \right\}
        \quad \text{(policy core)}
\end{align}
where $\theta_{\mathrm{policy}} > \theta_{\mathrm{norm}}$ is the elevated threshold
for positions eligible for actual policy implementation.
\end{definition}

\begin{proposition}[Overton Window as Psychological Attractor]
\label{prop:overton-attractor}
The Overton Window $\mathcal{O}(t)$ satisfies the axioms of a psychological attractor
(Definition~\ref{def:psych-attractor}):
\begin{enumerate}[label=(\roman*)]
    \item \textbf{Invariance}: The discourse fractal $\Frac_{\mathcal{D}}$ preserves
          acceptability: $\Frac_{\mathcal{D}}(\mathcal{O}) \subseteq \mathcal{O}$
    \item \textbf{Attraction}: Positions near the boundary tend toward the interior
          or exterior under social dynamics
    \item \textbf{Minimality}: The window is the smallest stable configuration supporting
          coherent public discourse
\end{enumerate}
\end{proposition}

\begin{proof}
\textbf{(i) Invariance}: Social reinforcement dynamics operate via the fractal
$\Frac_{\mathcal{D}} = \langle 1 : 5 : 7 \rangle$ (awareness $\to$ reception $\to$ rhythm).
Positions within $\mathcal{O}$ receive attention ($\noe{1}$), are absorbed into consensus
($\noe{5}$), and become habituated ($\noe{7}$). This self-reinforcing loop maintains the
window: $\Frac_{\mathcal{D}}(\mathcal{O}) \subseteq \mathcal{O}$.

\textbf{(ii) Attraction}: Define the basin as $U = \{x \mid \mu_{\mathrm{accept}}(x) >
\theta_{\mathrm{norm}} - \delta\}$. Positions in $U \setminus \mathcal{O}$ experience
normalization pressure: repeated exposure shifts $\mu_{\mathrm{accept}}$ upward until
$x \in \mathcal{O}$ or $x$ is expelled to $\mathcal{D} \setminus U$.

\textbf{(iii) Minimality}: If $\mathcal{O}' \subset \mathcal{O}$ satisfied (i) and (ii),
then positions in $\mathcal{O} \setminus \mathcal{O}'$ would lack stable discourse
grounding, contradicting the attractor definition.
\end{proof}

%% ============================================================================
%% SUBSECTION: MEASURE OVERTON WINDOW
%% ============================================================================
\subsection{Measuring the Overton Window: Discourse Fractal Analysis}

\begin{definition}[Discourse Fractal]
\label{def:discourse-fractal}
The \textbf{discourse fractal} $\Frac_{\mathcal{D}}$ governing acceptability dynamics is:
\[
    \Frac_{\mathcal{D}} = \langle 1 : 7 : 5 : 2 \rangle
\]
representing the cycle: Awareness ($\noe{1}$) $\to$ Rhythm/Habituation ($\noe{7}$)
$\to$ Reception/Absorption ($\noe{5}$) $\to$ Attraction/Acceptance ($\noe{2}$).
\end{definition}

\begin{definition}[Measure Overton Window Function]
\label{def:measure-overton}
The function $\mathtt{measure\_overton\_window} : \mathcal{D} \times \mathcal{T} \to
\mathcal{P}(\mathcal{D}) \times \mathbb{R}^3$ computes the current window state:
\[
\boxed{
    \mathtt{measure\_overton\_window}(\mathcal{D}, t) :=
    \left( \mathcal{O}(t), \; \mu_{\mathrm{center}}(t), \; w(t), \; \sigma(t) \right)
}
\]
where:
\begin{align}
    \mathcal{O}(t) &= \left\{ x \in \mathcal{D} \;\Big|\;
        \mu_{\mathrm{accept}}(x, t) \geq \theta_{\mathrm{norm}}(t) \right\} \\
    \mu_{\mathrm{center}}(t) &= \frac{1}{|\mathcal{O}(t)|} \sum_{x \in \mathcal{O}(t)}
        \mathrm{pos}(x) \quad \text{(window centroid)} \\
    w(t) &= \max_{x \in \mathcal{O}(t)} \mathrm{pos}(x) -
        \min_{x \in \mathcal{O}(t)} \mathrm{pos}(x) \quad \text{(window width)} \\
    \sigma(t) &= \sqrt{\frac{1}{|\mathcal{O}(t)|} \sum_{x \in \mathcal{O}(t)}
        (\mathrm{pos}(x) - \mu_{\mathrm{center}}(t))^2} \quad \text{(window variance)}
\end{align}
and $\mathrm{pos} : \mathcal{D} \to \mathbb{R}$ is a one-dimensional projection
mapping positions to a ``left-right'' or ``conservative-progressive'' spectrum.
\end{definition}

\begin{algorithm}[Discourse Fractal Analysis]
\label{alg:discourse-analysis}
\textbf{Input}: Discourse corpus $C$, position set $\mathcal{D}$, time $t$

\textbf{Output}: Window measurement $(\mathcal{O}, \mu_c, w, \sigma)$

\begin{enumerate}
    \item \textbf{Extract mentions}: For each $x \in \mathcal{D}$, count
          $n_x = |\{c \in C \mid x \text{ mentioned in } c\}|$
    \item \textbf{Compute sentiment}: $s_x = \frac{1}{n_x} \sum_{\text{mentions}}
          \mathrm{sentiment}(\text{context})$
    \item \textbf{Calculate acceptability}:
          $\mu_{\mathrm{accept}}(x) = \frac{n_x}{\max_y n_y} \cdot \frac{1 + s_x}{2}$
    \item \textbf{Determine threshold}:
          $\theta_{\mathrm{norm}} = \mathrm{median}(\{\mu_{\mathrm{accept}}(x)\}_{x \in \mathcal{D}})$
    \item \textbf{Construct window}: $\mathcal{O} = \{x \mid \mu_{\mathrm{accept}}(x)
          \geq \theta_{\mathrm{norm}}\}$
    \item \textbf{Compute statistics}: Calculate $\mu_c$, $w$, $\sigma$ per Definition~\ref{def:measure-overton}
\end{enumerate}
\end{algorithm}

%% ============================================================================
%% SUBSECTION: SHIFT WINDOW
%% ============================================================================
\subsection{Window Shifting via Transfinite Sequences}

\begin{definition}[Shift Fractal]
\label{def:shift-fractal}
The \textbf{shift fractal} $\Frac_{\mathrm{shift}}$ is a transfinite sequence of
acceptable positions:
\[
\boxed{
    \Frac_{\mathrm{shift}} := \left( x_\alpha \right)_{\alpha < \lambda}
    \quad \text{where } x_\alpha \in \mathcal{D} \text{ and }
    d(x_\alpha, x_{\alpha+1}) \leq \delta_{\max}
}
\]
for some limit ordinal $\lambda$ and maximum step size $\delta_{\max}$.

The sequence satisfies the \textbf{coherence condition}:
\[
    \forall \alpha < \beta < \lambda :
    \mu_{\mathrm{accept}}(x_\alpha, t_\alpha) \geq \theta_{\mathrm{norm}}(t_\alpha)
    \implies \mu_{\mathrm{accept}}(x_\beta, t_\beta) \geq \theta_{\mathrm{norm}}(t_\beta) - \epsilon
\]
Each step maintains acceptability within tolerance $\epsilon$ of the threshold.
\end{definition}

\begin{definition}[Shift Window Function]
\label{def:shift-window}
The function $\mathtt{shift\_window} : \mathcal{O} \times \mathcal{D} \times
\mathbb{R}_{>0} \to \mathcal{O}'$ implements a single transfinite increment:
\[
\boxed{
    \mathtt{shift\_window}(\mathcal{O}, x_{\mathrm{target}}, \delta) :=
    \Frac_{\mathcal{D}}^n\left( \mathcal{O} \cup \{x_{\mathrm{edge}}\} \right)
}
\]
where:
\begin{align}
    x_{\mathrm{edge}} &= \argmin_{x \in \mathcal{D} \setminus \mathcal{O}}
        \left\{ d(x, \partial\mathcal{O}) \;\Big|\;
        d(x, x_{\mathrm{target}}) < d(\mathcal{O}, x_{\mathrm{target}}) \right\} \\
    n &= \min\left\{ k \;\Big|\;
        \mu_{\mathrm{accept}}\left(\Frac_{\mathcal{D}}^k(x_{\mathrm{edge}})\right)
        \geq \theta_{\mathrm{norm}} \right\}
\end{align}
That is: select the nearest outside position that moves toward the target, then
iterate the discourse fractal until that position becomes normalized.
\end{definition}

\begin{theorem}[Shift Convergence]
\label{thm:shift-convergence}
Let $\mathcal{O}_0 = \mathcal{O}(t_0)$ be the initial window and
$x_{\mathrm{target}} \in \mathcal{D}$ the target position. Under the shift fractal
$\Frac_{\mathrm{shift}}$, the sequence of windows $(\mathcal{O}_\alpha)_{\alpha < \lambda}$
converges:
\[
    \lim_{\alpha \to \lambda} \mathcal{O}_\alpha = \mathcal{O}^*
    \quad \text{where } x_{\mathrm{target}} \in \mathrm{Core}(\mathcal{O}^*)
\]
provided the Collage Theorem conditions (Theorem~\ref{thm:collage}) are satisfied:
\[
    d_H(\mathcal{O}^*, \FractalAttractor_{\Frac_{\mathrm{shift}}}) \leq
    \frac{1}{1 - r_{\mathrm{shift}}} \cdot d_H(\mathcal{O}^*,
    H_{\Frac_{\mathrm{shift}}}(\mathcal{O}^*))
\]
\end{theorem}

\begin{proof}
The shift fractal $\Frac_{\mathrm{shift}}$ defines a contractive IFS on window space
with contractivity ratio $r_{\mathrm{shift}} = 1 - \delta_{\max} / \mathrm{diam}(\mathcal{D})$.
By the Banach Fixed Point Theorem (underlying the Collage Theorem), the sequence
$\mathcal{O}_\alpha = H_{\Frac_{\mathrm{shift}}}^\alpha(\mathcal{O}_0)$ converges to
the unique attractor $\FractalAttractor_{\Frac_{\mathrm{shift}}}$.

The collage error bound ensures that by designing $\Frac_{\mathrm{shift}}$ such that
$H_{\Frac_{\mathrm{shift}}}(\mathcal{O}^*)$ closely overlaps the target window
containing $x_{\mathrm{target}}$, the limiting attractor $\FractalAttractor$
approximates this target.
\end{proof}

%% ============================================================================
%% SUBSECTION: DESIGN WINDOW SHIFT
%% ============================================================================
\subsection{Designing Shift Strategies: Constraint Optimization}

\begin{definition}[Shift Strategy]
\label{def:shift-strategy}
A \textbf{shift strategy} $\mathcal{S}$ is a tuple:
\[
    \mathcal{S} = \left( x_{\mathrm{target}}, \; \delta_{\max}, \;
    \theta_{\mathrm{plausible}}, \; T_{\mathrm{max}}, \; \mathcal{C} \right)
\]
where:
\begin{itemize}
    \item $x_{\mathrm{target}} \in \mathcal{D}$: the desired final position
    \item $\delta_{\max} \in \mathbb{R}_{>0}$: maximum step size per iteration
    \item $\theta_{\mathrm{plausible}} \in [0,1]$: plausible deniability threshold
    \item $T_{\mathrm{max}} \in \mathbb{R}_{>0}$: maximum time horizon
    \item $\mathcal{C}$: constraint set (social, legal, resource bounds)
\end{itemize}
\end{definition}

\begin{definition}[Design Window Shift Function]
\label{def:design-window-shift}
The function $\mathtt{design\_window\_shift}$ solves the constrained optimization:
\[
\boxed{
    \mathtt{design\_window\_shift}(\mathcal{O}_0, x_{\mathrm{target}}, \mathcal{C}) :=
    \argmin_{\Frac_{\mathrm{shift}}} \left\{
        |\Frac_{\mathrm{shift}}| \;\Big|\;
        x_{\mathrm{target}} \in \mathcal{O}_{\infty} \text{ and }
        \Frac_{\mathrm{shift}} \models \mathcal{C}
    \right\}
}
\]
where $|\Frac_{\mathrm{shift}}|$ is the length (ordinal) of the shift sequence and
$\Frac_{\mathrm{shift}} \models \mathcal{C}$ denotes constraint satisfaction.
\end{definition}

\begin{definition}[Constraint Set for Window Shifting]
\label{def:shift-constraints}
The standard constraint set $\mathcal{C}_{\mathrm{std}}$ includes:
\begin{align}
    \mathcal{C}_{\mathrm{step}} &: \forall \alpha : d(x_\alpha, x_{\alpha+1}) \leq \delta_{\max}
        \quad \text{(step size bound)} \\
    \mathcal{C}_{\mathrm{plausible}} &: \forall \alpha :
        P(\text{``shift detected''} \mid x_\alpha) \leq \theta_{\mathrm{plausible}}
        \quad \text{(plausible deniability)} \\
    \mathcal{C}_{\mathrm{time}} &: \lambda \cdot \Delta t \leq T_{\mathrm{max}}
        \quad \text{(time horizon)} \\
    \mathcal{C}_{\mathrm{resource}} &: \sum_\alpha \mathrm{cost}(x_\alpha \to x_{\alpha+1})
        \leq B \quad \text{(budget constraint)} \\
    \mathcal{C}_{\mathrm{reversibility}} &: \forall \alpha :
        \exists \beta > \alpha : d(x_\beta, x_0) < d(x_\alpha, x_0)
        \quad \text{(escape path)}
\end{align}
\end{definition}

\begin{algorithm}[Optimal Shift Design via Collage Minimization]
\label{alg:shift-design}
\textbf{Input}: Current window $\mathcal{O}_0$, target $x_{\mathrm{target}}$,
constraints $\mathcal{C}$

\textbf{Output}: Optimal shift fractal $\Frac^*_{\mathrm{shift}}$

\begin{enumerate}
    \item \textbf{Construct target window}:
          $\mathcal{O}^* = \mathcal{O}_0 \cup B_\epsilon(x_{\mathrm{target}})$
          where $B_\epsilon$ is the $\epsilon$-ball
    \item \textbf{Initialize}: $S \gets \varnothing$, $\mathcal{O}_{\mathrm{curr}}
          \gets \mathcal{O}_0$
    \item \textbf{While} $x_{\mathrm{target}} \notin \mathcal{O}_{\mathrm{curr}}$:
        \begin{enumerate}
            \item Find $x_{\mathrm{edge}} = \argmin_{x \in \partial^+\mathcal{O}_{\mathrm{curr}}}
                  d(x, x_{\mathrm{target}})$
            \item \textbf{If} $d(x_{\mathrm{edge}}, \mathcal{O}_{\mathrm{curr}}) \leq \delta_{\max}$
                  \textbf{and} constraints $\mathcal{C}$ satisfied:
                \begin{enumerate}
                    \item $S \gets S \cup \{x_{\mathrm{edge}}\}$
                    \item $\mathcal{O}_{\mathrm{curr}} \gets \mathtt{shift\_window}
                          (\mathcal{O}_{\mathrm{curr}}, x_{\mathrm{target}}, \delta_{\max})$
                \end{enumerate}
            \item \textbf{Else}: Reduce $\delta_{\max}$ or relax constraints
        \end{enumerate}
    \item \textbf{Return} $\Frac^*_{\mathrm{shift}} = (x_\alpha)_{\alpha \leq |S|}$
          where $x_\alpha$ is the $\alpha$-th element of $S$
\end{enumerate}

\textbf{Optimality}: By the Collage Theorem, minimizing collage error
$d_H(\mathcal{O}^*, H_S(\mathcal{O}^*))$ minimizes the number of steps required.
\end{algorithm}

%% ============================================================================
%% SUBSECTION: NORMALIZATION DYNAMICS
%% ============================================================================
\subsection{Normalization Dynamics: Acceptance Threshold Evolution}

\begin{definition}[Acceptance Threshold Dynamics]
\label{def:acceptance-dynamics}
The \textbf{acceptance threshold} $\theta_{\mathrm{norm}}(t)$ evolves according to:
\[
\boxed{
    \frac{d\theta_{\mathrm{norm}}}{dt} =
    -\gamma \cdot (\theta_{\mathrm{norm}} - \bar{\mu}(t)) +
    \eta \cdot \sigma_{\mu}(t) \cdot \xi(t)
}
\]
where:
\begin{itemize}
    \item $\bar{\mu}(t) = \frac{1}{|\mathcal{D}|} \sum_{x \in \mathcal{D}}
          \mu_{\mathrm{accept}}(x, t)$: mean acceptability
    \item $\sigma_{\mu}(t)$: variance of acceptability distribution
    \item $\gamma > 0$: regression rate toward mean
    \item $\eta > 0$: noise sensitivity
    \item $\xi(t)$: stochastic perturbation (social shocks, events)
\end{itemize}
\end{definition}

\begin{definition}[Wait for Normalization Function]
\label{def:wait-normalization}
The function $\mathtt{wait\_for\_normalization}$ computes the time required for a
position to achieve acceptability:
\[
\boxed{
    \mathtt{wait\_for\_normalization}(x, \theta_{\mathrm{target}}) :=
    \inf\left\{ t > 0 \;\Big|\; \mu_{\mathrm{accept}}(x, t) \geq \theta_{\mathrm{target}} \right\}
}
\]
\end{definition}

\begin{theorem}[Normalization Time Bound]
\label{thm:normalization-time}
For a position $x$ at initial acceptability $\mu_0 = \mu_{\mathrm{accept}}(x, 0) <
\theta_{\mathrm{norm}}$, the expected normalization time satisfies:
\[
    \mathbb{E}\left[ \mathtt{wait\_for\_normalization}(x, \theta_{\mathrm{norm}}) \right]
    \leq \frac{1}{\alpha} \cdot \ln\left( \frac{\theta_{\mathrm{norm}} - \mu_0}{\epsilon} \right)
\]
where $\alpha$ is the exposure rate and $\epsilon$ is the convergence tolerance.
\end{theorem}

\begin{proof}
Under repeated exposure via the discourse fractal $\Frac_{\mathcal{D}}$, acceptability
evolves as:
\[
    \frac{d\mu_{\mathrm{accept}}}{dt} = \alpha \cdot (\theta_{\mathrm{target}} -
    \mu_{\mathrm{accept}}) + \beta \cdot \mathrm{exposure}(x, t)
\]
This is an exponential approach with rate $\alpha$. Solving:
\[
    \mu_{\mathrm{accept}}(t) = \theta_{\mathrm{target}} -
    (\theta_{\mathrm{target}} - \mu_0) \cdot e^{-\alpha t}
\]
Setting $\mu_{\mathrm{accept}}(T) = \theta_{\mathrm{norm}} - \epsilon$ and solving for $T$
yields the bound.
\end{proof}

\begin{definition}[Completion Criterion]
\label{def:completion-criterion}
The window shift is \textbf{complete} when:
\[
\boxed{
    \mathtt{shift\_complete}(\mathcal{O}, x_{\mathrm{target}}) :=
    \left[ x_{\mathrm{target}} \in \mathrm{Core}(\mathcal{O}) \right] \land
    \left[ \frac{d\mu_{\mathrm{accept}}(x_{\mathrm{target}})}{dt} \geq 0 \right]
}
\]
That is, the target is in the policy core and its acceptability is stable or increasing.
\end{definition}

%% ============================================================================
%% SECTION: CANONICAL EQUATIONS
%% ============================================================================
\section{Canonical Equations}
\label{sec:overton-canonical}

\begin{tcolorbox}[colback=gray!5!white,colframe=gray!75!black,title=\textbf{Canonical Overton Window Equations},breakable]

\textbf{Equation 1: Window Measurement}
\begin{equation}
\boxed{
    \mathcal{O}(t) = \left\{ x \in \mathcal{D} \;\Big|\;
    \mu_{\mathrm{accept}}(x, t) \geq \theta_{\mathrm{norm}}(t) \right\}
}
\tag{OW-1}
\end{equation}

\textbf{Equation 2: Shift Trajectory}
\begin{equation}
\boxed{
    \mathcal{O}_{\alpha+1} = \Frac_{\mathcal{D}}^{n_\alpha}\left(
    \mathcal{O}_\alpha \cup \{x_{\alpha+1}\} \right)
    \quad \text{where } n_\alpha =
    \min\{k \mid \mu_{\mathrm{accept}}(\Frac_{\mathcal{D}}^k(x_{\alpha+1}))
    \geq \theta_{\mathrm{norm}}\}
}
\tag{OW-2}
\end{equation}

\textbf{Equation 3: Step Size Constraint}
\begin{equation}
\boxed{
    d(x_\alpha, x_{\alpha+1}) \leq \delta_{\max} \cdot
    \left(1 + \frac{\mu_{\mathrm{accept}}(x_\alpha) - \theta_{\mathrm{norm}}}
    {\sigma_{\mu}}\right)
}
\tag{OW-3}
\end{equation}

\textbf{Equation 4: Normalization Dynamics}
\begin{equation}
\boxed{
    \frac{d\mu_{\mathrm{accept}}(x)}{dt} =
    \alpha \cdot \mathrm{exposure}(x, t) \cdot
    \left( \theta_{\mathrm{ceiling}} - \mu_{\mathrm{accept}}(x) \right) -
    \beta \cdot \left( \mu_{\mathrm{accept}}(x) - \mu_{\mathrm{floor}} \right) \cdot
    \mathbf{1}_{[\mathrm{no\ exposure}]}
}
\tag{OW-4}
\end{equation}

\textbf{Equation 5: Completion Criterion}
\begin{equation}
\boxed{
    \mathtt{complete} \iff
    x_{\mathrm{target}} \in \mathrm{Core}(\mathcal{O}) \land
    d_H\left(\mathcal{O}, \FractalAttractor_{\Frac_{\mathrm{shift}}}\right) < \epsilon
}
\tag{OW-5}
\end{equation}

\end{tcolorbox}

%% ============================================================================
%% SECTION: PLAIN ENGLISH EXPLANATION
%% ============================================================================
\section{Plain English Explanation}
\label{sec:overton-plain}

\begin{tcolorbox}[colback=green!5!white,colframe=green!65!black,title=\textbf{Plain English: What is the Overton Window?},breakable]

\textbf{The Simple Version:}
The Overton Window is the range of ideas that society considers ``acceptable to discuss''
at any given time. Ideas inside the window can be debated openly; ideas outside are
dismissed as extreme, crazy, or unthinkable. The window is not fixed---it shifts
over time as society changes.

\textbf{The ``Boiling Frog'' Metaphor:}
\begin{itemize}
    \item Imagine a frog in a pot of water. If you heat the water slowly enough,
          the frog doesn't notice the gradual temperature change and stays in the pot.
    \item Similarly, societal acceptance shifts gradually. What seems radical today
          becomes debatable tomorrow and policy the day after---if the shift is
          slow enough that people adapt without noticing.
    \item Each small step seems reasonable compared to the previous position.
    \item Looking back over time, the total shift can be dramatic.
\end{itemize}

\textbf{The ``Moving the Goalposts'' Metaphor:}
\begin{itemize}
    \item Imagine a soccer goal that moves slightly each time a shot is taken.
    \item The \textbf{window} is where the goalposts currently stand.
    \item \textbf{Shifting} is moving the goalposts a little bit at a time.
    \item \textbf{Normalization} is when players stop noticing the new position.
    \item \textbf{Plausible deniability} is moving slowly enough that no single move
          looks suspicious.
\end{itemize}

\textbf{Technical Term $\leftrightarrow$ Metaphor Mapping:}
\begin{center}
\footnotesize
\begin{tabularx}{\textwidth}{>{\raggedright\arraybackslash}X|>{\raggedright\arraybackslash}X|>{\raggedright\arraybackslash}X}
\toprule
\textbf{Technical Term} & \textbf{Metaphor (Goalposts)} & \textbf{Metaphor (Frog)} \\
\midrule
$\mathcal{O}$ (current window) & Where the goalposts stand now & Current water temperature \\
$x_{\mathrm{target}}$ (target position) & Where you want the goal & Target temperature \\
$\delta_{\max}$ (max step size) & How far you move per shift & How much you heat per minute \\
$\theta_{\mathrm{norm}}$ (threshold) & What counts as ``in bounds'' & Comfort threshold \\
$\mu_{\mathrm{accept}}$ (acceptability) & How normal a position seems & Frog's comfort level \\
$\Frac_{\mathrm{shift}}$ (shift fractal) & Sequence of small moves & Heating schedule \\
Plausible deniability & ``It was always there'' & ``Water was always warm'' \\
Normalization time & Time until new normal & Time until frog adapts \\
\bottomrule
\end{tabularx}
\end{center}

\textbf{Why This Matters:}
Understanding the Overton Window helps explain how dramatic social changes occur
through incremental steps. The same mathematical process describes:
\begin{itemize}
    \item How civil rights expanded (beneficial shift)
    \item How privacy norms eroded with technology (contested shift)
    \item How extremist ideas enter mainstream discourse (concerning shift)
\end{itemize}
The mathematics is neutral---it describes the \textit{mechanism}, not the morality.

\end{tcolorbox}

%% ============================================================================
%% SECTION: TECHNICAL TERM MAPPING TABLE
%% ============================================================================
\section{Technical Term to Metaphor Mapping}
\label{sec:overton-mapping}

\begin{center}
\small
\renewcommand{\arraystretch}{1.4}
\begin{tabularx}{\textwidth}{|>{\raggedright\arraybackslash}p{3.5cm}|>{\raggedright\arraybackslash}X|>{\raggedright\arraybackslash}X|}
\hline
\textbf{Technical Term} & \textbf{Plain English Meaning} & \textbf{Real-World Example} \\
\hline\hline
\texttt{current\_window} $(\mathcal{O})$ &
The range of ideas currently considered acceptable in mainstream discussion &
In 2000: ``Marriage is between man and woman'' was mainstream; same-sex marriage was outside window \\
\hline
\texttt{target\_window} $(\mathcal{O}^*)$ &
The desired range of acceptable ideas after the shift completes &
In 2015: Same-sex marriage is legal and mainstream; opposition is increasingly outside window \\
\hline
\texttt{shift\_strategy} $(\mathcal{S})$ &
The plan for moving from current to target window &
Civil rights strategy: Legal challenges $\to$ Media representation $\to$ State laws $\to$ Federal ruling \\
\hline
\texttt{max\_step\_size} $(\delta_{\max})$ &
How big a change can happen in one step without backlash &
``Civil unions'' as intermediate step; ``Domestic partnerships'' as earlier step \\
\hline
\texttt{plausible\_deniability} $(\theta_{\mathrm{plausible}})$ &
Ability to claim no deliberate manipulation is occurring &
``We're just asking questions'' or ``This is just about equality'' \\
\hline
\texttt{normalization\_threshold} $(\theta_{\mathrm{norm}})$ &
How accepted an idea must be to enter mainstream discussion &
When polls show $>40\%$ support, media begins ``both sides'' coverage \\
\hline
\texttt{edge\_case\_articles} &
Media pieces that test the boundary by presenting ``extreme'' positions &
Op-eds by fringe figures that ``spark conversation'' about previously unthinkable positions \\
\hline
\texttt{discourse\_fractal} $(\Frac_{\mathcal{D}})$ &
The repetitive process by which ideas become normalized &
Exposure $\to$ Familiarity $\to$ Acceptance $\to$ Advocacy cycle \\
\hline
\texttt{wait\_for\_normalization} &
Time required for a new position to become mainstream &
Gay rights: Stonewall (1969) to Marriage Equality (2015) = 46 years \\
\hline
\texttt{collage\_error} &
How far the current shift is from the target &
Distance between current policy debates and ultimate goal \\
\hline
\end{tabularx}
\end{center}

%% ============================================================================
%% SECTION: REAL-LIFE SCENARIOS
%% ============================================================================
\section{Real-Life Scenarios}
\label{sec:overton-scenarios}

\begin{tcolorbox}[colback=yellow!5!white,colframe=yellow!65!black,title=\textbf{Real-Life Scenarios: Overton Window Shifts},breakable]

\textbf{Scenario 1: Civil Rights Expansion (1950s--2015)}

The movement for marriage equality represents a successful, multi-generational
window shift.

\textit{Historical sequence:}
\begin{enumerate}
    \item \textbf{1950s}: Homosexuality classified as mental illness; discussion taboo
    \item \textbf{1969}: Stonewall riots---public visibility begins
    \item \textbf{1970s--80s}: ``Gay rights'' enters vocabulary; legal battles begin
    \item \textbf{1990s}: ``Domestic partnerships'' as compromise position
    \item \textbf{2000s}: ``Civil unions'' gain acceptance; first state marriages
    \item \textbf{2015}: Marriage equality federally recognized; opposition becomes fringe
\end{enumerate}

\textit{Technical mapping:}
\begin{itemize}
    \item $\mathcal{O}_{1950} = \{\text{heterosexual marriage only}\}$
    \item $x_{\mathrm{target}} = \text{marriage equality}$
    \item $\Frac_{\mathrm{shift}} = \langle x_1, x_2, \ldots, x_n \rangle$ where
          $x_i$ = ``decriminalization,'' ``domestic partnership,'' ``civil union,''
          ``state marriage,'' ``federal marriage''
    \item $\delta_{\max} \approx$ one legal status change per decade
    \item $\mathtt{wait\_for\_normalization}(x_{\mathrm{target}}) \approx 46$ years
    \item Completion: $x_{\mathrm{target}} \in \mathrm{Core}(\mathcal{O}_{2015})$ via
          \textit{Obergefell v. Hodges}
\end{itemize}

\textbf{Canonical equation application:}
\[
    \mathcal{O}_{2015} = \Frac_{\mathcal{D}}^{n}\left( \mathcal{O}_{1950} \cup
    \{x_1, x_2, \ldots, x_n\} \right)
    \quad \text{where } n \approx 65 \text{ years of normalization cycles}
\]

\hrulefill

\textbf{Scenario 2: Technology Privacy Norms (1990s--Present)}

The erosion of privacy expectations demonstrates window shifting driven by
technological and commercial forces.

\textit{Historical sequence:}
\begin{enumerate}
    \item \textbf{1990s}: Strong privacy expectations; encryption debates
    \item \textbf{2001}: Post-9/11 surveillance normalization begins
    \item \textbf{2004--2010}: Social media makes sharing personal data routine
    \item \textbf{2013}: Snowden revelations---temporary backlash, then acceptance resumes
    \item \textbf{2020s}: Facial recognition, location tracking accepted as normal
\end{enumerate}

\textit{Technical mapping:}
\begin{itemize}
    \item $\mathcal{O}_{1995} = \{\text{strong personal privacy},
          \text{government needs warrant}\}$
    \item $x_{\mathrm{target}} = \text{pervasive data collection accepted}$
    \item $\delta_{\max}$ = each new platform/service normalizes one more data type
    \item Normalization mechanism: Convenience trade-off (``free'' services for data)
    \item $\theta_{\mathrm{plausible}} =$ ``Users agreed to Terms of Service''
\end{itemize}

\textbf{Canonical equation application:}
\[
    \frac{d\mu_{\mathrm{accept}}(\text{surveillance})}{dt} =
    \alpha \cdot \mathrm{convenience} - \beta \cdot \mathrm{privacy\_awareness}
\]
where $\alpha \gg \beta$ in practice, leading to steady erosion.

\hrulefill

\textbf{Scenario 3: Economic Policy---Taxation Acceptability (1950s--2020s)}

Changes in acceptable tax rates demonstrate bidirectional window movement.

\textit{Historical sequence:}
\begin{enumerate}
    \item \textbf{1950s}: Top marginal rate 91\%---accepted as normal
    \item \textbf{1960s--70s}: Gradual reduction; 70\% becomes new ceiling
    \item \textbf{1980s}: Reagan cuts to 28\%; supply-side economics enters mainstream
    \item \textbf{1990s--2000s}: Any rate above 40\% labeled ``socialist''
    \item \textbf{2010s--20s}: Wealth taxes, higher rates re-enter discourse
\end{enumerate}

\textit{Technical mapping:}
\begin{itemize}
    \item Window shift direction: First rightward (lower rates), then partial leftward return
    \item $\Frac_{\mathrm{shift,1}} = \langle 91\%, 70\%, 50\%, 28\%, 35\% \rangle$ (1950--2000)
    \item $\Frac_{\mathrm{shift,2}} = \langle 35\%, 39.6\%, 50\%?, \ldots \rangle$ (2010--present)
    \item $\delta_{\max} \approx 10$--15 percentage points per political generation
\end{itemize}

\textbf{Canonical equation application:}
\[
    \mathcal{O}_{\mathrm{tax}}(t) = \left\{ r \in [0, 100] \;\Big|\;
    |r - \bar{r}(t)| \leq w(t)/2 \right\}
\]
where $\bar{r}(t)$ is the politically ``moderate'' rate and $w(t)$ is window width.

\hrulefill

\textbf{Scenario 4: Medical Ethics---Procedure Acceptability}

The evolution of end-of-life care options illustrates ethical window shifting.

\textit{Historical sequence:}
\begin{enumerate}
    \item \textbf{Pre-1970s}: Life extension at all costs; euthanasia unthinkable
    \item \textbf{1970s}: ``Right to die'' enters discourse; Quinlan case
    \item \textbf{1990s}: DNR orders normalized; hospice care expands
    \item \textbf{1997}: Oregon Death with Dignity Act---physician-assisted death legal
    \item \textbf{2010s--20s}: Multiple states adopt similar laws; ``MAID'' terminology
\end{enumerate}

\textit{Technical mapping:}
\begin{itemize}
    \item $\mathcal{O}_{1960} = \{\text{preserve life unconditionally}\}$
    \item $x_{\mathrm{target}} = \text{patient autonomy over end-of-life}$
    \item Intermediate positions: DNR, palliative sedation, hospice, aid-in-dying
    \item $\Frac_{\mathrm{shift}} = \langle 3 : 1 : 5 : 2 \rangle$ (reject old $\to$
          awareness $\to$ absorb new $\to$ accept)
    \item $\theta_{\mathrm{plausible}} =$ ``Patient's choice,'' ``Compassion,''
          ``Dignity''
\end{itemize}

\textbf{Canonical equation application:}
\[
    \mu_{\mathrm{accept}}(\text{MAID}, t) = \mu_0 +
    \int_0^t \alpha \cdot \mathrm{media\_exposure}(\tau) \cdot
    (1 - \mu_{\mathrm{accept}}(\tau)) \, d\tau
\]
Saturation model where exposure drives acceptance toward ceiling.

\end{tcolorbox}

%% ============================================================================
%% SECTION: CONSTRUCTION RULES
%% ============================================================================
\section{Overton Window Equation Construction Rules}
\label{sec:overton-construction}

\begin{tcolorbox}[colback=yellow!5!white,colframe=yellow!75!black,title=\textbf{Construction Algorithm},breakable]

To construct Overton Window equations for a given scenario:

\textbf{Step 1: Define the Discourse Space}
\begin{itemize}
    \item Identify the domain: politics, ethics, technology, economics, etc.
    \item Map positions to a spectrum: $\mathrm{pos} : \mathcal{D} \to \mathbb{R}$
    \item Identify current ``center'' and ``edges''
\end{itemize}

\textbf{Step 2: Measure Current Window}
\begin{itemize}
    \item Apply $\mathtt{measure\_overton\_window}$ to determine $\mathcal{O}(t)$
    \item Compute $\mu_{\mathrm{center}}$, $w$, $\sigma$
    \item Identify boundary positions $\partial^{\pm}\mathcal{O}$
\end{itemize}

\textbf{Step 3: Specify Target}
\begin{itemize}
    \item Define $x_{\mathrm{target}}$: the position to be normalized
    \item Compute distance: $d(x_{\mathrm{target}}, \mathcal{O})$
    \item Estimate required shift magnitude
\end{itemize}

\textbf{Step 4: Design Shift Fractal}
\begin{itemize}
    \item Apply Collage Theorem to minimize $d_H(\mathcal{O}^*, H_S(\mathcal{O}^*))$
    \item Construct intermediate positions: $\Frac_{\mathrm{shift}} =
          \langle x_1, x_2, \ldots, x_n \rangle$
    \item Verify each step satisfies $d(x_i, x_{i+1}) \leq \delta_{\max}$
\end{itemize}

\textbf{Step 5: Apply Constraints}
\begin{itemize}
    \item Check plausible deniability: no single step too large
    \item Verify time horizon: total shift time realistic
    \item Ensure resource bounds: each step has implementation path
\end{itemize}

\textbf{Step 6: Model Normalization}
\begin{itemize}
    \item Apply Equation (OW-4) to each intermediate position
    \item Compute $\mathtt{wait\_for\_normalization}(x_i)$ for each step
    \item Verify cumulative time within $T_{\max}$
\end{itemize}

\textbf{Step 7: Verify Completion}
\begin{itemize}
    \item Check $x_{\mathrm{target}} \in \mathrm{Core}(\mathcal{O}_{\mathrm{final}})$
    \item Verify stability: $d\mu_{\mathrm{accept}}/dt \geq 0$
    \item Confirm attractor convergence via Equation (OW-5)
\end{itemize}

\end{tcolorbox}

%% ============================================================================
%% SECTION: CLASSIFICATION TABLE
%% ============================================================================
\section{Overton Window Shift Classification}
\label{sec:overton-classification}

\begin{center}
\small
\renewcommand{\arraystretch}{1.3}
\begin{longtable}{p{2.2cm}|p{2cm}|p{2.5cm}|p{2cm}|p{3cm}}
\caption{Overton Window Shift Classifications} \\
\toprule
\textbf{Shift Type} & \textbf{Direction} & \textbf{Characteristic Fractal} & \textbf{Typical Duration} & \textbf{Historical Example} \\
\midrule
\endfirsthead
\multicolumn{5}{c}{\textit{Table continued from previous page}} \\
\toprule
\textbf{Shift Type} & \textbf{Direction} & \textbf{Characteristic Fractal} & \textbf{Typical Duration} & \textbf{Historical Example} \\
\midrule
\endhead
\midrule
\multicolumn{5}{r}{\textit{Continued on next page}} \\
\endfoot
\bottomrule
\endlastfoot
Rights Expansion & Progressive & $\langle 1:2:5:7 \rangle$ & 30--60 years & Civil rights, marriage equality \\
Privacy Erosion & Both ways & $\langle 4:5:7 \rangle$ & 10--20 years & Social media norms \\
Economic Shift & Oscillating & $\langle 7:2:3:7 \rangle$ & 20--40 years & Tax policy \\
Medical Ethics & Progressive & $\langle 3:1:5:2 \rangle$ & 20--40 years & End-of-life care \\
Technology Adoption & Progressive & $\langle 4:2:5:7 \rangle$ & 5--15 years & AI, genetic editing \\
Environmental Policy & Contested & $\langle 1:3:2:7 \rangle$ & 30--50 years & Climate policy \\
Security vs Liberty & Reactive & $\langle 8:3:5 \rangle$ & Event-driven & Post-crisis surveillance \\
Cultural Values & Both ways & $\langle 7:5:7 \rangle$ & Generation-length & Religious observance \\
\end{longtable}
\end{center}

%% ============================================================================
%% SECTION: MATHEMATICAL PROPERTIES
%% ============================================================================
\section{Mathematical Properties of Window Dynamics}
\label{sec:overton-properties}

\begin{theorem}[Window Monotonicity]
\label{thm:window-monotonicity}
Under sustained one-directional pressure, the Overton Window centroid shifts
monotonically:
\[
    \forall t_1 < t_2 : \mathrm{sign}\left(
        \frac{d\mu_{\mathrm{center}}}{dt} \right) = \mathrm{const}
    \implies \mu_{\mathrm{center}}(t_2) - \mu_{\mathrm{center}}(t_1) \text{ has same sign}
\]
\end{theorem}

\begin{theorem}[Width Oscillation]
\label{thm:width-oscillation}
The window width $w(t)$ oscillates with social tension:
\[
    w(t) = w_0 + A \cdot \cos(\omega t + \phi) + \epsilon(t)
\]
where high-tension periods compress the window (smaller $w$) and low-tension
periods expand it.
\end{theorem}

\begin{theorem}[Irreversibility Threshold]
\label{thm:irreversibility}
A window shift becomes effectively irreversible when:
\[
    \mu_{\mathrm{accept}}(x_{\mathrm{old}}) < \theta_{\mathrm{norm}} -
    3\sigma_{\mu}
\]
That is, the former ``normal'' position has moved more than three standard
deviations below the acceptance threshold.
\end{theorem}

\begin{corollary}[Generational Lock-In]
\label{cor:generational}
If a position $x$ achieves $\mu_{\mathrm{accept}}(x) \geq \theta_{\mathrm{policy}}$
and remains there for one generation ($\approx 25$ years), then
$\mu_{\mathrm{accept}}(x)$ becomes resistant to reversal:
\[
    \frac{d\mu_{\mathrm{accept}}(x)}{dt} \approx 0 \quad \text{for } t > t_0 + 25
\]
The position has entered the ``common sense'' attractor.
\end{corollary}

%% ============================================================================
%% SECTION: CONNECTION TO TKS CANON
%% ============================================================================
\section{Connection to TKS Canonical Structures}
\label{sec:overton-tks-connection}

\begin{proposition}[Overton Window as Noetic Attractor]
\label{prop:overton-noetic}
The Overton Window $\mathcal{O}$ corresponds to a B-world psychological attractor
in TKS:
\[
    \mathcal{O} \cong \PsychAttractor_{B,\mathrm{collective}}
\]
where the governing fractal is:
\[
    \Frac_{\mathcal{O}} = \langle 1 : 7 : 5 : 2 \rangle_{2b}
\]
applying the Wisdom sub-foundation ($\Foundations_{2b}$) in the intellectual world.
\end{proposition}

\begin{proposition}[Window Shift as RPM Process]
\label{prop:overton-rpm}
A designed window shift corresponds to an RPM monad computation:
\[
    \mathtt{shift\_window} : \RPM(\mathcal{O} \to \mathcal{O}')
\]
with:
\begin{itemize}
    \item Desire chain $D_n$: The target position $x_{\mathrm{target}}$
    \item Wisdom chain $W_n$: The shift strategy $\mathcal{S}$
    \item Power chain $P_n$: The implementation sequence $\Frac_{\mathrm{shift}}$
\end{itemize}
\end{proposition}

\begin{proposition}[Collage Theorem Application]
\label{prop:overton-collage}
The optimal window shift design is the TKS Collage Theorem (Theorem~\ref{thm:collage})
applied to discourse space:
\[
    d_H(\mathcal{O}^*, \FractalAttractor_{\Frac_{\mathrm{shift}}}) \leq
    \frac{1}{1 - r_{\mathrm{shift}}} \cdot d_H(\mathcal{O}^*,
    H_{\Frac_{\mathrm{shift}}}(\mathcal{O}^*))
\]
Minimizing the collage error yields the most efficient shift path.
\end{proposition}

%% ============================================================================
%% SECTION: SUMMARY
%% ============================================================================
\section{Chapter Summary}
\label{sec:overton-summary}

\begin{tcolorbox}[colback=blue!5!white,colframe=blue!75!black,title=\textbf{Key Results},breakable]

This chapter established the TKS formalization of Overton Window dynamics:

\begin{enumerate}
    \item \textbf{Definition}: The Overton Window is a psychological attractor basin
          in discourse space, characterized by acceptability measure $\mu_{\mathrm{accept}}$
          and normalization threshold $\theta_{\mathrm{norm}}$.

    \item \textbf{Measurement}: The $\mathtt{measure\_overton\_window}$ function
          extracts window parameters via discourse fractal analysis.

    \item \textbf{Shifting}: Window shifts occur via transfinite sequences of
          acceptable positions, governed by the shift fractal $\Frac_{\mathrm{shift}}$.

    \item \textbf{Design}: Optimal shift strategies are computed via Collage Theorem
          minimization subject to step-size, plausibility, and resource constraints.

    \item \textbf{Normalization}: Acceptance dynamics follow exponential approach
          models with exposure-dependent rates.

    \item \textbf{Completion}: A shift is complete when the target enters the
          policy core with stable or increasing acceptability.
\end{enumerate}

\textbf{Canonical Equations:}
\begin{itemize}
    \item (OW-1): Window measurement
    \item (OW-2): Shift trajectory
    \item (OW-3): Step size constraint
    \item (OW-4): Normalization dynamics
    \item (OW-5): Completion criterion
\end{itemize}

\end{tcolorbox}

\begin{tcolorbox}[colback=red!5!white,colframe=red!75!black,title=\textbf{Methodological Note}]
The mathematics presented in this chapter is \textbf{value-neutral}. The same equations
describe:
\begin{itemize}
    \item Progressive social movements expanding rights
    \item Corporate campaigns normalizing data collection
    \item Authoritarian regimes shifting discourse boundaries
    \item Organic cultural evolution over generations
\end{itemize}
The formalization provides tools for \textit{understanding} these dynamics, not for
\textit{endorsing} any particular direction of change. Ethical evaluation of specific
applications requires normative frameworks beyond the scope of TKS mathematics.
\end{tcolorbox}
